\vspace{-0.2cm}
\section{Discussion, Limitations, and Conclusion}
\vspace{-0.2cm}

In this work, we investigated how to imbue LLMs with self-correction behavior that enables them to correct their own responses on the fly. To accomplish this, we proposed \methodname{}, a multi-turn RL approach, and demonstrated through extensive evaluations \emph{that it is one of the first methods} to attain significantly positive intrinsic self-correction performance. To do so, we rigorously analyzed the behavior of various SFT approaches and identified failure modes in which the model learns a non-correcting strategy (e.g. learning to make no edits; behavior collapse) or falls prey to distribution shift. \methodname{} trains a self-correcting strategy by utilizing a two-stage design and reward shaping, both of which help preventing behavior collapse into not learning effective self-corrective behavior. 
\methodname{} has limitations that also provide avenues for future work. We did not train \methodname{} for more than one round of iterative self-correction due to infrastructural reasons, which means that subsequent rounds may not be as effective as the first. Future work should train with more than two attempts via RL, which is already a common and effective practice to obtain effective self-correction behavior over more than two rounds with SFT~\citep{snell2024scaling,qu2024recursive}. Unifying Stages I and II would also be interesting, since it would alleviate the limitation of running multiple runs. Finally, our results suggest that learning \emph{meta-strategies} (e.g., self-correction) might require going beyond standard LLM fine-tuning (Section~\ref{sec:sft}), and incorporate regularization (e.g., progress reward).


